\chapter{Elettrochimica}
 
L'elettrochimica studia le reazioni nelle quali avviene un trasferimento di elettroni fra le specie coinvolte, e il modo in cui tale trasferimento può essere sfruttato per produrre lavoro elettrico oppure, all'inverso, forzato dall'esterno per far avvenire trasformazioni che spontaneamente non avverrebbero.
Il filo conduttore dell'intero capitolo è che una reazione redox non è altro che una reazione con $\Delta G$ come tutte le altre, e che la novità sta soltanto nel fatto che il cammino degli elettroni può essere separato spazialmente dai reagenti.
 
\section{Reazioni di ossidoriduzione}
 
Si dice ossidazione\mkeyword{ossidazione} la perdita di elettroni da parte di una specie e riduzione\mkeyword{riduzione} il loro acquisto.

I due processi non possono avvenire separatamente, poiché gli elettroni non esistono liberi in soluzione: la specie che si ossida li cede a quella che si riduce, e il numero complessivo di elettroni ceduti deve eguagliare quello degli elettroni acquistati.

La specie che si riduce, cioè quella che sottrae elettroni all'altra, si chiama ossidante\mkeyword{ossidante}; quella che si ossida, e che quindi li fornisce, si chiama riducente\mkeyword{riducente}.
 
Per tenere il conto degli elettroni si introduce il numero di ossidazione\mkeyword{numero di ossidazione}, cioè la carica che un atomo avrebbe se tutti i legami che lo coinvolgono fossero completamente ionici, attribuendo la coppia di legame all'atomo più elettronegativo.
Non è una grandezza fisica, dal momento che la densità elettronica reale è distribuita con continuità e non a numeri interi, ma è un artificio contabile perfettamente adeguato allo scopo: una reazione è redox quando almeno un atomo cambia numero di ossidazione, e la variazione misura il numero di elettroni scambiati.
 
\section{La pila}
 
Se si immerge una lamina di zinco in una soluzione di solfato di rame, la reazione
\begin{equation}
    \mathrm{Zn}_{(s)} + \mathrm{Cu^{2+}}_{(aq)} \longrightarrow \mathrm{Zn^{2+}}_{(aq)} + \mathrm{Cu}_{(s)}
    \label{eq:daniell}
\end{equation}
avviene spontaneamente per contatto diretto: gli elettroni passano dallo zinco allo ione rame nel punto stesso in cui i due si toccano, e l'energia liberata si degrada integralmente in calore.\\

L'idea della pila\mkeyword{pila} consiste nel separare fisicamente le due semireazioni, costringendo gli elettroni a percorrere un circuito esterno per andare da una parte all'altra: la stessa reazione, con lo stesso $\Delta G$, produce allora lavoro elettrico anziché calore.

\newpage

La cella è costituita da due semielementi\mkeyword{semielemento}, ciascuno formato da un elettrodo metallico immerso nella soluzione di un suo sale.
Nel caso in esame lo zinco immerso in $\mathrm{Zn(NO_3)_2}$ costituisce l'anodo\mkeyword{anodo}, sede dell'ossidazione

\begin{equation*}
    \mathrm{Zn}_{(s)} \longrightarrow \mathrm{Zn^{2+}}_{(aq)} + 2e^- ,
\end{equation*}

mentre il rame immerso in $\mathrm{Cu(NO_3)_2}$ costituisce il catodo\mkeyword{catodo}, sede della riduzione

\begin{equation*}
    \mathrm{Cu^{2+}}_{(aq)} + 2e^- \longrightarrow \mathrm{Cu}_{(s)} .
\end{equation*}

I due semielementi sono collegati esternamente da un filo conduttore, nel quale scorrono gli elettroni, e internamente da un ponte salino\mkeyword{ponte salino} contenente un elettrolita indifferente, tipicamente $\mathrm{NaNO_3}$, i cui ioni non partecipano ad alcuna reazione ma migrano in modo da mantenere la neutralità elettrica delle due soluzioni.
Senza il ponte salino la soluzione anodica accumulerebbe carica positiva e quella catodica carica negativa, e il flusso di elettroni si arresterebbe quasi immediatamente.

\begin{figure}[htbp]
    \centering
    \begin{tikzpicture}[>={Latex[length=2mm]}, font=\small]
        % --- macro riutilizzabili -----------------------------------------
        % becher aperto in alto: #1 = ascissa del centro
        \def\becher#1{%
            \draw (#1-1.6,0.2) -- (#1-1.6,-2.6) -- (#1+1.6,-2.6) -- (#1+1.6,0.2);
            \draw (#1-1.6,-0.4) -- (#1+1.6,-0.4);
        }
        % elettrodo: #1 = ascissa, #2 = etichetta, #3 = segno, #4 = verso esterno
        \def\elettrodo#1#2#3#4{%
            \draw (#1-0.09,-2.0) rectangle (#1+0.09,1.0);
            \node at (#1+0.42*#4,0.70) {#3};
            \node at (#1+0.50*#4,-1.40) {#2};
        }
        % --- becher ed elettrodi ------------------------------------------
        \becher{-3.0}
        \becher{3.0}
        \elettrodo{-3.0}{$\mathrm{Zn}$}{$-$}{-1}
        \elettrodo{3.0}{$\mathrm{Cu}$}{$+$}{1}
        % --- circuito esterno con voltmetro -------------------------------
        \draw (-3.0,1.0) -- (-3.0,2.4) -- (-0.55,2.4);
        \draw (0.55,2.4) -- (3.0,2.4) -- (3.0,1.0);
        \draw (0,2.4) circle (0.55);
        \node at (0,2.4) {$V$};
        \draw[->] (-2.2,2.4) -- (-1.7,2.4);
        \draw[->] (1.7,2.4) -- (2.2,2.4);
        \node[anchor=south] at (-1.95,2.5) {$e^-$};
        \node[anchor=south] at (1.95,2.5) {$e^-$};
        % --- ponte salino --------------------------------------------------
        \draw[rounded corners=10pt]
            (-1.9,-1.2) -- (-1.9,0.9) -- (1.9,0.9) -- (1.9,-1.2);
        \draw[->] (-0.35,0.9) -- (-0.95,0.9);
        \draw[->] (0.35,0.9) -- (0.95,0.9);
        \node[anchor=south] at (-0.95,0.98) {$\mathrm{NO_3^-}$};
        \node[anchor=south] at (0.95,0.98) {$\mathrm{Na^+}$};
        % --- specie in soluzione -------------------------------------------
        \node at (-2.15,-2.15) {$\mathrm{Zn^{2+}}$};
        \node at (2.15,-2.15) {$\mathrm{Cu^{2+}}$};
        % --- semireazioni ---------------------------------------------------
        \node[anchor=north] at (-3.0,-2.8)
            {$\mathrm{Zn}_{(s)} \rightarrow \mathrm{Zn^{2+}}_{(aq)} + 2e^-$};
        \node[anchor=north] at (3.0,-2.8)
            {$\mathrm{Cu^{2+}}_{(aq)} + 2e^- \rightarrow \mathrm{Cu}_{(s)}$};
    \end{tikzpicture}
    \caption{Schema della pila Daniell}
    \label{fig:pila_daniell}
\end{figure}

Conviene fissare alcuni punti che sono fonte costante di confusione.
L'anodo produce elettroni e il catodo li consuma; gli elettroni non attraversano mai la soluzione, ma soltanto il circuito esterno, e vanno sempre dall'anodo al catodo; nella soluzione la corrente è trasportata dagli ioni, non dagli elettroni.
Ne segue che in una pila l'anodo è il polo negativo e il catodo il polo positivo, mentre nell'elettrolisi, come si vedrà, i segni sono invertiti pur restando invariata la definizione di anodo e catodo in termini di ossidazione e riduzione.
 
La cella si rappresenta in forma compatta scrivendo l'anodo a sinistra e il catodo a destra, una barra singola per ogni interfaccia fra fasi diverse e una barra doppia per il ponte salino:

\begin{equation}
    \mathrm{Zn}_{(s)} \ | \ \mathrm{Zn^{2+}}_{(aq)} \ || \ \mathrm{Cu^{2+}}_{(aq)} \ | \ \mathrm{Cu}_{(s)} .
    \label{eq:notazione_cella}
\end{equation}
 
\newpage

\section{Potenziali di elettrodo}
 
La differenza di potenziale che si misura fra i due elettrodi a circuito aperto si chiama forza elettromotrice\mkeyword{forza elettromotrice} della pila e si indica con $E$.
Quando tutte le specie in soluzione hanno concentrazione $1 \ M$, i gas eventualmente presenti hanno pressione di $1 \ atm$ e la temperatura è di $25 \ ^\circ C$, si parla di potenziale standard\mkeyword{potenziale standard} e si scrive $E^\circ$.
 
Il potenziale di un singolo semielemento non è misurabile in modo assoluto, perché per misurare una differenza di potenziale occorrono sempre due elettrodi.
Si sceglie allora un riferimento convenzionale, l'elettrodo standard a idrogeno\mkeyword{elettrodo standard a idrogeno}, nel quale la semireazione

\begin{equation*}
    2\mathrm{H^+}_{(aq)} + 2e^- \longrightarrow \mathrm{H_2}_{(g)}
\end{equation*}

avviene in condizioni standard su una superficie di platino, e si pone per definizione $E^\circ = 0 \ V$.
Tutti gli altri potenziali sono misurati accoppiando il semielemento in esame a questo riferimento, e si tabulano per convenzione come potenziali di riduzione\mkeyword{potenziale di riduzione}, cioè scrivendo sempre la semireazione nel verso della riduzione.
La forza elettromotrice standard della pila è allora

\begin{equation}
    E^\circ_{cella} = E^\circ_{red}(\text{catodo}) - E^\circ_{red}(\text{anodo}) ,
    \label{eq:fem_standard}
\end{equation}

e la pila funziona spontaneamente nel verso scritto quando $E^\circ_{cella} > 0$, cioè quando il potenziale di riduzione del catodo è maggiore di quello dell'anodo.

La scala dei potenziali standard ordina le specie secondo il loro potere ossidante.
In cima si trova il fluoro, con $E^\circ = +2{,}87 \ V$, l'ossidante più forte fra quelli comuni; in fondo il litio, con $E^\circ = -3{,}04 \ V$, il riducente più forte.
La posizione dell'idrogeno a $0 \ V$ separa i metalli che riescono a spostarlo dagli acidi, cioè quelli con potenziale negativo come zinco, ferro e magnesio, da quelli che non ci riescono, come rame, argento e oro, i quali hanno potenziale positivo e per essere disciolti richiedono acidi che siano anche ossidanti.
 
\section{Il legame con la termodinamica}
 
Il lavoro elettrico massimo che una pila può compiere è il prodotto della carica trasportata per la differenza di potenziale, e coincide con la variazione di energia libera cambiata di segno, poiché $\Delta G$ misura appunto il lavoro utile massimo a $T$ e $P$ costanti.
Se $n$ è il numero di moli di elettroni scambiate per mole di reazione e $F = 96\,500 \ C/mol$ è la costante di Faraday\mkeyword{costante di Faraday}, cioè la carica di una mole di elettroni, si ottiene
\begin{equation}
    \Delta G = -nFE ,
    \qquad\qquad
    \Delta G^\circ = -nFE^\circ .
    \label{eq:deltag_nfe}
\end{equation}
La \eqref{eq:deltag_nfe} traduce immediatamente il criterio di spontaneità: la reazione è spontanea, $\Delta G < 0$, se e solo se $E > 0$.
Il segno della forza elettromotrice sostituisce dunque il segno dell'energia libera, e la misura di una tensione con un voltmetro è a tutti gli effetti una misura calorimetrica indiretta.
 
Da qui segue anche la regola per combinare semireazioni, che è la fonte di errore più comune.
Poiché $\Delta G$ è additivo mentre $E$ non lo è, per sommare due semireazioni occorre passare per le energie libere.
Volendo per esempio ricavare $E^\circ$ della coppia $\mathrm{Fe^{3+}}/\mathrm{Fe}$ a partire da
\[
    \mathrm{Fe^{3+}} + e^- \longrightarrow \mathrm{Fe^{2+}} , \quad E^\circ = +0{,}77 \ V ,
    \qquad
    \mathrm{Fe^{2+}} + 2e^- \longrightarrow \mathrm{Fe} , \quad E^\circ = -0{,}44 \ V ,
\]
si sommano le due energie libere e si divide per i tre elettroni complessivi,
\begin{equation}
    E^\circ = \frac{n_1 E^\circ_1 + n_2 E^\circ_2}{n_1 + n_2}
    = \frac{1 \cdot 0{,}77 + 2 \cdot (-0{,}44)}{3} = -0{,}04 \ V ,
    \label{eq:combinazione_potenziali}
\end{equation}
che è una media pesata sul numero di elettroni, non una somma.
Sommare direttamente i due potenziali avrebbe dato $+0{,}33 \ V$, un valore privo di senso fisico.
 
Infine, poiché $\Delta G^\circ = -RT \ln K$, confrontando con la \eqref{eq:deltag_nfe} si ottiene il legame fra potenziale standard e costante di equilibrio,
\begin{equation}
    \ln K = \frac{nFE^\circ}{RT} ,
    \qquad\text{ovvero}\qquad
    \log K = \frac{n E^\circ}{0{,}0592} ,
    \label{eq:logk}
\end{equation}
dove il valore numerico $0{,}0592 \ V$ è $RT\ln 10/F$ calcolato a $298 \ K$.
Una differenza di potenziale di appena $0{,}0592 \ V$ per elettrone corrisponde quindi a un fattore dieci sulla costante di equilibrio, il che spiega perché potenziali di pochi decimi di volt bastino a rendere una reazione praticamente quantitativa.
 
\section{Equazione di Nernst}
 
Una pila continua a funzionare finché $E > 0$, e si arresta quando $E = 0$: a quel punto, per la \eqref{eq:deltag_nfe}, si ha anche $\Delta G = 0$, cioè la reazione ha raggiunto l'equilibrio e la pila è scarica\mkeyword{pila scarica}.

Conviene sottolineare che una pila esaurita non è una pila «vuota»: i reagenti sono ancora presenti, ma la miscela si trova all'equilibrio e da essa non è più estraibile alcun lavoro utile.
Il punto in cui questo avviene è determinato dalle concentrazioni delle specie che partecipano alla reazione, poiché la capacità ossidante di una specie $\mathrm{Ox}$ cresce con la sua concentrazione, e con essa cresce $E$.
 
La relazione si ricava direttamente dalla termodinamica.
Partendo da $\Delta G = \Delta G^\circ + RT \ln Q$ e sostituendo la \eqref{eq:deltag_nfe} in entrambi i membri si ottiene $-nFE = -nFE^\circ + RT \ln Q$, cioè
\begin{equation}
    E = E^\circ - \frac{RT}{nF} \ln Q .
    \label{eq:nernst_cella}
\end{equation}
Applicata a una singola semireazione di riduzione $\mathrm{Ox} + ne^- \rightarrow \mathrm{Red}$, per la quale $Q_{sem} = [\mathrm{Red}]/[\mathrm{Ox}]$, la \eqref{eq:nernst_cella} si scrive più comodamente invertendo l'argomento del logaritmo, il che cambia il segno del termine e dà l'equazione di Nernst\mkeyword{equazione di Nernst}
\begin{equation}
    E_{red} = E^\circ_{red} + \frac{RT}{nF} \ln \frac{[\mathrm{Ox}]}{[\mathrm{Red}]} .
    \label{eq:nernst}
\end{equation}
Con $[\mathrm{Ox}]$ e $[\mathrm{Red}]$ si intendono i prodotti delle concentrazioni di tutte le specie che compaiono rispettivamente a sinistra e a destra della semireazione, ciascuna elevata al proprio coefficiente stechiometrico; le fasi condensate pure vi compaiono con attività unitaria e i gas con la loro pressione parziale.
 
A $25 \ ^\circ C$, includendo nella costante anche il fattore di conversione al logaritmo decimale, si ha $RT \ln 10 / F = 0{,}0592 \ V$, e la \eqref{eq:nernst} assume la forma d'uso corrente
\begin{equation}
    E_{red} = E^\circ_{red} + \frac{0{,}0592}{n} \log \frac{[\mathrm{Ox}]}{[\mathrm{Red}]} .
    \label{eq:nernst_log}
\end{equation}
La lettura qualitativa è immediata: più l'ossidante è concentrato, più il potenziale della coppia è elevato; man mano che la pila lavora l'ossidante si consuma e il riducente si accumula, cosicché $E$ diminuisce fino ad annullarsi.
 
Qualche esempio chiarisce come si costruisce il termine logaritmico.
Per la coppia $\mathrm{Zn^{2+}}/\mathrm{Zn}$ il metallo è una fase condensata pura e non compare, sicché
\[
    E_{red} = -0{,}76 + \frac{0{,}0592}{2} \log [\mathrm{Zn^{2+}}] .
\]
Per la coppia $\mathrm{H^+}/\mathrm{H_2}$ compare la pressione parziale dell'idrogeno al denominatore,
\[
    E_{red} = 0 + \frac{0{,}0592}{2} \log \frac{[\mathrm{H^+}]^2}{P_{\mathrm{H_2}}} ,
\]
che a pressione unitaria si riduce a $E_{red} = -0{,}0592 \ \mathrm{pH}$.
Per una coppia in cui entrambe le forme sono in soluzione, come $\mathrm{Fe^{3+}}/\mathrm{Fe^{2+}}$, si ha semplicemente $E_{red} = E^\circ_{red} + 0{,}0592 \log([\mathrm{Fe^{3+}}]/[\mathrm{Fe^{2+}}])$.
Il caso più laborioso è quello del permanganato, dove la semireazione
\[
    \mathrm{MnO_4^-} + 5e^- + 8\mathrm{H^+} \longrightarrow \mathrm{Mn^{2+}} + 4\mathrm{H_2O}
\]
comporta il passaggio del manganese dal numero di ossidazione $+7$ a $+2$, mentre ossigeno e idrogeno non cambiano il proprio, da cui i cinque elettroni; l'acqua è il solvente e non compare, mentre gli ioni $\mathrm{H^+}$ entrano con l'esponente otto,
\[
    E_{red} = E^\circ_{red} + \frac{0{,}0592}{5} \log \frac{[\mathrm{MnO_4^-}][\mathrm{H^+}]^8}{[\mathrm{Mn^{2+}}]} .
\]
La dipendenza dall'ottava potenza di $[\mathrm{H^+}]$ spiega perché il permanganato sia un ossidante efficace soltanto in ambiente fortemente acido: ogni unità di $\mathrm{pH}$ in più abbassa il potenziale di quasi un decimo di volt.
 
Nota infine la forza elettromotrice della cella si ottiene, come sempre, dalla differenza fra il potenziale della coppia che si riduce e quello della coppia che si ossida, entrambi scritti come potenziali di riduzione e ciascuno calcolato con la \eqref{eq:nernst_log} alle concentrazioni effettive.

\section{Elettrolisi}
 
L'elettrolisi\mkeyword{elettrolisi} è il processo inverso: applicando dall'esterno una differenza di potenziale maggiore in modulo di quella della pila corrispondente si costringe la reazione ad avvenire nel verso non spontaneo, convertendo lavoro elettrico in energia chimica.
Le definizioni di anodo e catodo restano quelle di sempre, l'anodo sede dell'ossidazione e il catodo sede della riduzione, ma i segni si invertono: nella cella elettrolitica l'anodo è collegato al polo positivo del generatore e il catodo a quello negativo.
 